\documentclass[letterpaper]{article} % DO NOT CHANGE THIS
\usepackage[submission]{aaai2027}  % DO NOT CHANGE THIS
% The serif, sans-serif, and monospaced fonts are loaded automatically by
% aaai2027.sty (newtxtext, helvet, courier). DO NOT add \usepackage{times},
% \usepackage{helvet}, \usepackage{courier}, or any other font package.
\usepackage[hyphens]{url}  % DO NOT CHANGE THIS
\usepackage{graphicx} % DO NOT CHANGE THIS
\urlstyle{rm} % DO NOT CHANGE THIS
\def\UrlFont{\rm}  % DO NOT CHANGE THIS
\usepackage{natbib}  % DO NOT CHANGE THIS AND DO NOT ADD ANY OPTIONS TO IT
\usepackage{caption} % DO NOT CHANGE THIS AND DO NOT ADD ANY OPTIONS TO IT
\frenchspacing  % DO NOT CHANGE THIS
%
% These are recommended to typeset algorithms but not required. See the subsubsection on algorithms. Remove them if you don't have algorithms in your paper.
\usepackage{algorithm}
\usepackage{algorithmic}
\usepackage{multirow}
\usepackage{enumitem}
\usepackage[figuresright]{rotating}
\usepackage[table]{xcolor}
\usepackage{booktabs}
\usepackage{array}
\usepackage{makecell}
\usepackage{multirow}
\usepackage{pdflscape}
\usepackage{amssymb}
%
% These are recommended to typeset listings but not required. See the subsubsection on listing. Remove this block if you don't have listings in your paper.
\usepackage{newfloat}
\usepackage{listings}
\DeclareCaptionStyle{ruled}{labelfont=normalfont,labelsep=colon,strut=off} % DO NOT CHANGE THIS
\lstset{%
	basicstyle={\footnotesize\ttfamily},% footnotesize acceptable for monospace
	numbers=left,numberstyle=\footnotesize,xleftmargin=2em,% show line numbers, remove this entire line if you don't want the numbers.
	aboveskip=0pt,belowskip=0pt,%
	showstringspaces=false,tabsize=2,breaklines=true}
\floatstyle{ruled}
\newfloat{listing}{tb}{lst}{}
\floatname{listing}{Listing}

%
% Recommended for better-looking tables
\usepackage{booktabs}

%
% Keep the \pdfinfo as shown here. There's no need
% for you to add the /Title and /Author tags.
\pdfinfo{
/TemplateVersion (2027.1)
}

% DISALLOWED PACKAGES
% \usepackage{authblk} -- This package is specifically forbidden
% \usepackage{balance} -- This package is specifically forbidden
% \usepackage{CJK} -- This package is specifically forbidden
% \usepackage{float} -- This package is specifically forbidden
% \usepackage{flushend} -- This package is specifically forbidden
% \usepackage{fullpage} -- This package is specifically forbidden
% \usepackage{geometry} -- This package is specifically forbidden
% \usepackage{hyperref} -- This package is specifically forbidden
% \usepackage{navigator} -- This package is specifically forbidden
% (or any other package that embeds links such as navigator or hyperref)
% \usepackage{indentfirst} -- This package is specifically forbidden
% \usepackage{layout} -- This package is specifically forbidden
% \usepackage{multicol} -- This package is specifically forbidden
% \usepackage{nameref} -- This package is specifically forbidden
% \usepackage{savetrees} -- This package is specifically forbidden
% \usepackage{setspace} -- This package is specifically forbidden
% \usepackage{stfloats} -- This package is specifically forbidden
% \usepackage{tabu} -- This package is specifically forbidden
% \usepackage{titlesec} -- This package is specifically forbidden
% \usepackage{tocbibind} -- This package is specifically forbidden
% \usepackage{ulem} -- This package is specifically forbidden
% \usepackage{wrapfig} -- This package is specifically forbidden
% DISALLOWED COMMANDS
% \nocopyright -- Your paper will not be published if you use this command
% \addtolength -- This command may not be used
% \balance -- This command may not be used
% \baselinestretch -- Your paper will not be published if you use this command
% \clearpage -- No page breaks of any kind may be used for the final version of your paper
% \columnsep -- This command may not be used
% \newpage -- No page breaks of any kind may be used for the final version of your paper
% \pagebreak -- No page breaks of any kind may be used for the final version of your paper
% \pagestyle -- This command may not be used
% \tiny -- This is not an acceptable font size.
% \vspace{- -- No negative value may be used in proximity of a caption, figure, table, section, subsection, subsubsection, or reference
% \vskip{- -- No negative value may be used to alter spacing above or below a caption, figure, table, section, subsection, subsubsection, or reference

\setcounter{secnumdepth}{0} %May be changed to 1 or 2 if section numbers are desired.

% The file aaai2027.sty is the style file for AAAI Press
% proceedings, working notes, and technical reports.
%

% Title

% Your title must be in mixed case, not sentence case.
% That means all verbs (including short verbs like be, is, using,and go),
% nouns, adverbs, adjectives should be capitalized, including both words in hyphenated terms, while
% articles, conjunctions, and prepositions are lower case unless they
% directly follow a colon or long dash
\title{Sci-VLA: Agentic VLA Inference Plugin for Long-Horizon Tasks in Scientific Experiments}


\author {
    % Authors
    First Author Name\textsuperscript{\rm 1,\rm 2}\equalcontrib,
    Second Author Name\textsuperscript{\rm 2}\equalcontrib,
    Third Author Name\textsuperscript{\rm 1}\corresponding
}
\affiliations {
    % Affiliations
    \textsuperscript{\rm 1}Affiliation 1\\
    \textsuperscript{\rm 2}Affiliation 2\\
    firstAuthor@affiliation1.com, secondAuthor@affilation2.com, thirdAuthor@affiliation1.com
}



\begin{document}

\maketitle

\begin{abstract}
Robotic laboratories play a critical role in autonomous scientific discovery by enabling scalable, continuous experimental execution. Recent vision–language–action (VLA) models offer a promising foundation for robotic laboratories. However, scientific experiments typically involve long-horizon tasks composed of multiple atomic tasks, posing a fundamental challenge to existing VLA models. While VLA models fine-tuned for scientific tasks can often execute \textit{atomic} experimental actions seen during training, they may fail to perform \textit{composite} tasks formed by reordering and composing these known atomic actions. This limitation arises from a distributional mismatch between training-time atomic tasks and inference-time composite tasks, which prevents VLA models from executing necessary transitional action between atomic tasks. To address this challenge, we propose Sci, an Agentic VLA Inference Plugin for Long-Horizon Tasks in Scientific Experiments. Sci introduces an LLM-based agentic inference mechanism that inserts transitional action between every two consecutive atomic tasks. By performing explicit transition inference and generating transitional robotic action code, Sci guides VLA models through missing transitional steps, improving the execution coherence of composite scientific workflows without additional VLA training. This inference-only intervention is data-efficient and training-free, while introducing additional LLM inference cost. We build 3D assets of scientific instruments and common scientific operating scenes within an existing simulation environment. In these scenes, we observe that Sci improves the independent success rate of atomic tasks during inference. Furthermore, we present a preliminary real-world demonstration of transferring Sci from simulation to a laboratory setting.
\end{abstract}

% \begin{teaserfigure}
%   \includegraphics[width=\textwidth]{Figures/figure0-teaserfigure.pdf}
%   %\vspace{-15px}
%   \caption{Vision-language-action (VLA) models often suffer from the state gap issue when inferring the open and long-horizon task in scientific scenarios. This paper proposes an agentic VLA inference plugin to generate transitional actions, thus bridging the gaps among composite tasks.
%   We verify its effectiveness on multiple 3-/5-/8-step tasks in the digital-twin system.}
%   \vspace{5px}
%   % \Description{Enjoying the baseball game from the third-base
%   % seats. Ichiro Suzuki preparing to bat.}
%   \label{fig0:teaser}
% \end{teaserfigure}

% Uncomment the following to link to your code, datasets, an extended version or similar.
% You must keep this block between (not within) the abstract and the main body of the paper.
% Make sure that you do not de-anonymize yourself with these links.
% \begin{links}
%     \link{Code}{https://aaai.org/example/code}
%     \link{Datasets}{https://aaai.org/example/datasets}
%     \link{Extended version}{https://aaai.org/example/extended-version}
% \end{links}


\begin{figure*}[!t]
    \centering
    \includegraphics[width=\linewidth]{Figures/intro.pdf}
    \caption{
    Given a VLA model fine-tuned for atomic tasks in scientific experiments,
    it may fail to complete the composite task due to the state gap issue.
    In this paper,
    we propose a LLM-based agentic inference plugin to generate transitional robotic execution code, improving execution of the long-horizon tasks.
    }
    \label{fig:motivation}
\end{figure*}


\section{Introduction}

In recent years, the AI4Science paradigm~\cite{wang2023scientific}
has evolved from single-model approaches~\cite{jumper2021highly,bi2023accurate}
toward autonomous scientific discovery~\cite{novikov2025alphaevolve,volk2023alphaflow,dai2025adaptive,koscher2023autonomous}.
Emerging research aims to directly embed artificial intelligence into scientific workflows to systematically enhance the entire research lifecycle~\cite{koscher2023autonomous,szymanski2023autonomous,boiko2023autonomous}.
In such a context, a robotic laboratory plays a critical role in enabling autonomous discovery~\cite{burger2020mobile,szymanski2023autonomous,coley2019robotic}.
It creates experimental environments with high-precision control, strong reproducibility, and large-scale execution capabilities, substantially accelerating scientific progress while reducing labor and time costs.
Moreover, robots capable of continuous, unattended operation significantly reduce researchers’ exposure to hazardous experimental conditions.
For example, mobile robotic chemists have demonstrated the feasibility of treating robots as human scientists~\cite {burger2020mobile}.
And autonomous laboratories have successfully discovered correct synthesis routes for 41 theoretically predicted materials through large-scale automated material synthesis experiments~\cite{szymanski2023autonomous}.



Current robotic laboratories typically follow the automation paradigm based on manually programmed, fixed experimental procedures
~\cite{coley2020autonomous}.
This choice is largely motivated by the high-precision requirements of scientific experiments.
For instance, in high-throughput chemical screening~\cite{macarron2011impact},
a single experimental cycle typically consists of sample preparation and reaction execution, for which fixed programs for robotic arms provide stability guarantees.
However, such rigid workflows limit the robot’s ability to reconfigure task sequences or respond to unforeseen situations, making them ineffective in open-ended experimental settings~\cite{ochiai2025automating}.
As an example, the parameters of a thermal cycler—such as cycle numbers, temperatures, and durations—must be flexibly adjusted according to experimental objectives, templates, and primers (the scientific instrument in Fig.~\ref{fig:motivation}), and thus cannot be pre-encoded as a single robot trajectory.
Researchers have to spend substantial time repeatedly reconfiguring robot parameters in the laboratory.


Compared with embodied intelligence approaches~\cite{an2025embodied} based on reinforcement learning~\cite{singh2022reinforcement},
recent vision–language–action (VLA) models~\cite{black2024pi_0,kim2025openvla} with strong generalization capabilities are better suited to open-ended robotic laboratories.
A typical VLA system integrates a vision–language model (VLM)~\cite{zhang2024vision,radford2021learning} with an action module (typically Transformer~\cite{vaswani2017attention}), enabling it to interpret natural language instructions, perceive the current visual scene, and generate low-level robotic control commands.
Unlike fixed-program systems, this paradigm allows robots to understand diverse natural language inputs and rapidly switch operational modes, making it potentially promising for open-ended laboratory scenarios.

In general, training data for VLA are difficult to acquire~\cite{mandlekar2023mimicgen}, and training or fine-tuning such models is computationally intensive~\cite{lin2024data}.
As the set of atomic experimental procedures are generally fixed and finite,
VLA models can be trained or fine-tuned using data collected only for \textit{atomic} tasks.
Then, to support diverse experimental protocols, these atomic procedures can be dynamically composed and reordered, resulting in the classic long-horizon challenge in robotics~\cite{fan2025long,gupta2020relay}.
Interestingly, we observe a long-horizon phenomenon of VLA for scientific experiments:
\textit{
a VLA model fine-tuned for scientific tasks can often execute atomic tasks seen during training, yet fails to correctly perform composite tasks constructed from these known atomic actions.}
As illustrated in Fig.~\ref{fig:motivation} (a), a fine-tuned VLA model can independently execute \texttt{open-thermal-cycler} (Task C.1) and \texttt{place-PCR-plate} (Task C.2), but fails to correctly execute the composite task \texttt{open–place}.
This limitation stems from the distributional mismatch between training-time data and inference-time task, which prevents VLA from executing the necessary transitional operations between atomic Task C.1 and Task C.2.

To address this challenge, we propose Sci, an Agentic VLA Inference Plugin for Long-Horizon Tasks in Scientific Experiments.
Given a VLA model fine-tuned for scientific experiments—i.e., one that has learned atomic tasks—we introduce an LLM-based agentic inference mechanism that inserts transitional action between every two consecutive atomic tasks.
The agent estimates the state gap, performs transition inference, and generates transitional robotic execution code corresponding to the transition, improving the execution coherence of composite scientific workflows.
This inference-only intervention eliminates the need for additional VLA training and is data-efficient, although it introduces additional LLM inference time during execution.
We design multiple long-horizon scientific operation tasks in a simulation environment, and then we test our method on these long-horizon tasks. Experimental results indicate that Sci improves the independent success rate of atomic tasks across various VLA models and improves the robot's execution coherence between continuous atomic tasks.



\section{Related Work}
\subsection{Robotic Laboratories}

Robotics labs primarily utilize automated processes involving robotic arms, wheeled robots, and drive shafts to replace human labor in laboratory experiments~\cite{boyd2002robotic}.
For example, self-driving thin-film laboratories have been developed for autonomous materials research~\cite{macleod2020self}. HTRobot~\cite{zhao2021discovery} can assemble approximately 1400 perovskite film arrangements.
A-Lab~\cite{szymanski2023autonomous} uses literature-trained language models and thermodynamically active sensing to propose a solid-state synthesis path.
A mobile robot
% chemists
chemist~\cite{burger2020mobile} automatically explores the ten-dimensional photocatalytic parameter space to find catalysts with performance nearly six times that of the initial formulation.
Adam~\cite{king2009automation} is the first scientist to autonomously discover the function of new genes in yeast through fully automated experiments and hypothesis testing.



\subsection{Robot Manipulation}
\textbf{Reinforcement Learning-based Methods.}
Reinforcement learning~\cite{singh2022reinforcement}
is a decision-making approach that interacts with the environment and uses rewards as the learning objective.
It has demonstrated success in certain real-world scenarios. For example, a reinforcement-learning robot that can make hot dogs with dual arms~\cite{li2019formal}. A human-in-the-loop reinforcement learning system that achieved strong performance on a wide range of dexterous manipulation tasks~\cite{luo2025precise}. A unified reinforcement learning–based control policy for the whole body to achieve effective shuttlecock tracking and striking~\cite{ma2025learning}.
However, robots based on reinforcement learning still find it difficult to be deployed in real open-world scenarios~\cite{ghosh2021generalization}. For example, in scientific experimental operations, continuous interactive learning in a real environment may lead to safety issues. Transferring knowledge from offline reinforcement learning to real-world scenarios is quite challenging~\cite{dulac2021challenges}.

\noindent
\textbf{Vision-Language-Action Models.}
Compared with the reinforcement learning paradigm, the presence of VLM
\cite{zhang2024vision,radford2021learning}
significantly improves the model's generalization performance in highly open-ended scenarios.
Given discrete time steps $t = \{1, \dots, T\}$,
visual $o_t^v$, language input $o^l \in \mathcal{O}_l$,
action space $a_t \in \mathcal{A}$, and combined information defined as: $\mathcal{H}_t = \{o^l,o_t^v \}$. The goal of VLA is to learn a strategy $\pi: \mathcal{H}_t \rightarrow \mathcal{A}$. Currently, there are two mainstream modeling approaches for VLA:
autoregression-based methods~\cite{zitkovich2023rt,kim2025openvla}
and diffusion-based methods~\cite{black2024pi_0,liurdt}.

Autoregression-based VLAs use combined information $\mathcal{H}_t$ to predict discrete action vectors. First, they map multimodal context to token sequences: visual observation token $z_t^{v} = f_v(o_t^{v}) \in \mathbb{R}^{N_v \times d}$ and language token $z^{l} = f_l(o^{l}) \in \mathbb{R}^{N_l \times d}$. Then, the policy function is implemented with a decoder-only transformer:
\begin{equation}
    p(z_t^a \mid \mathcal{H}_t) = \mathrm{Transformer}_\theta(\big[ z^{l}, z_1^{v}, \dots, z_{t-1}^{v}, z_t^{v} \big])
\end{equation}
where
the action token $z_t^a$ is detokenized into discrete action $a_t$:
\begin{equation}
    a_t = \mathrm{detokenize}(z_t^a).
\end{equation}
For example, RT-2~\cite{zitkovich2023rt} is an early end-to-end autoregression-based model.
Subsequently, OpenVLA~\cite{kim2025openvla} replaces the VLM base with an open-source large model, further accelerating VLA model development. $\pi_0\_fast$~\cite{pertsch2025fast} is an autoregressive VLA model that uses a fast tokenizer, enhancing training speed.

The diffusion-based method primarily follows the stepwise denoising paradigm of diffusion models, enabling VLA to predict continuous actions. Define an action chunk with a time length of $H$ and a starting time step $t$: $a_{t:t+H}$. Diffusion-based VLA implementing conditional probability distribution $p(a_{t:t+H} \mid \mathcal{H}_t, \tilde{a}_{t:t+H})$, where $\tilde{a}_{t:t+H} \sim \mathcal{N}(0, I)$ is a noisy action chunk. The final continuous action vector needs to be obtained through multiple denoising steps.
For example, diffusion policy~\cite{chi2025diffusion} uses the diffusion model idea to generate action trajectories. RDT~\cite{liurdt} uses more dual-arm robot training data and a larger transformer architecture, significantly improving zero-shot generalization performance.
$\pi_0$~\cite{black2024pi_0} and $\pi_{0.5}$~\cite{black2025pi_} use flow matching algorithms to gradually denoise actions.
They use the lightweight VLM: PaliGemma~\cite{beyer2024paligemma} as the base model. After fine-tuning on large-scale robot demonstration data, $\pi_0$ and $\pi_{0.5}$ achieve extremely strong generalization performance and can adapt to various household tasks.

\begin{figure*}[!t]
  \centering
  \includegraphics[width=\linewidth]{Figures/avip.pdf}
  \caption{The illustration of the Sci pipeline. Includes two modules: transitional action generation and transitional action insertion. The plugin and VLA will alternate in executing the sequence of atomic tasks.}
  \label{figs:avip}
\end{figure*}

\noindent
\textbf{Long-Horizon Challenges.}
Long-horizon tasks are complex tasks that span many time steps, typically formed by concatenating multiple atomic tasks.
Solving a long-horizon task involves decomposing it into independent atomic tasks and optimizing each one. In reinforcement learning, hierarchical reinforcement learning is primarily used to decompose and solve long-horizon tasks~\cite{dietterich2000hierarchical,konidaris2009skill}.
However, task decomposition can lead to issues such as error accumulation because it does not account for the connections and dependencies among atomic tasks. This is a challenge called skill chaining~\cite{konidaris2009skill} that needs to be addressed by hierarchical reinforcement learning methods. For example, a reinforcement learning-based dressing robot that aligns the state distributions of pre- and post-tasks to address the skill chaining problem~\cite{clegg2018learning}. Sequential dexterity is a reinforcement learning system that progressively finetunes the sub-policies to enhance the chaining success rate~\cite{chen2023sequential}.

VLA models leverage the task decomposition capabilities of VLM backbones or incorporate atomic task decomposition during training, which improve the model's generalization to long-horizon tasks. For example, $\pi_{0.5}$~\cite{black2025pi_} adds task decomposition data pairs during VLM pretraining.
LohoVLA~\cite{yang2025lohovla} generates token information for atomic tasks before generating action tokens. LongVLA~\cite{fan2025long} allows the model to focus on different visual information during different phases of movement and interaction, reducing the interference of redundant information during the task transition periods.
However, most of this work focuses on task decomposition and dataset processing, requiring training to achieve performance on long-horizon tasks.


% \clearpage
% \newpage

\section{Sci: Agentic VLA Inference Plugin}
We design Sci, an inference plugin that assists VLA execution without changing the VLA weights. Specifically, when an atomic task sequence arrives, Sci generates and inserts transitional action between every two consecutive atomic tasks, rather than only after a detected failure. A VLA model equipped with this plugin is denoted as Sci-\emph{model}, such as Sci-$\pi_0$.

% \begin{figure*}[!t]
%   \centering
%   \includegraphics[width=\linewidth]{Figures/3dasset.pdf}
%   %\vspace{-5px}
%   \caption{Some scientific simulated objects we used and their corresponding real object. Top: 3D simulated assets. Bottom: real objects. (a) Pipette. (b) Pipette rack. (c) Centrifuge tube rack. (d) Centrifuge tube plate. (e) 3D printer.}
%   \label{figs:3d1}
% \end{figure*}

\subsection{Inference Atomic Tasks Sequence}
In the Sci inference process, the generation module generates transitional action code after receiving the next atomic task $i+1$ prompt.
Then the insertion module interrupts VLA's inference when the pre-atomic-task $i$ finishes.
The transitional action is generated and inserted into the action sequence before the VLA starts atomic task $i+1$.
Finally, the VLA resumes inference for the subsequent task (See in Fig.~\ref{figs:avip}).
The two modules operate iteratively. This cycle continues until the entire long-horizon task is finished.

\textbf{Transitional Action Generation Module.}
When VLA's inference stops, the current information is obtained, including the current state $curr_s$, the current camera views $o^v_t$, and the description of the next atomic task $o^l_{i+1}$. Then, the generation module retrieves a possible target state $target_s$ from the training data $\mathcal{D}$ by the target task description $o_{i+1}^l$.
Specifically, we first extract a description set $\mathcal{P}=\{p_1, \dots, p_K\}$ containing all descriptions of atomic tasks, where $K$ stands for the number of training tasks. The retrieval module matches the next-task prompt to a training atomic-task prompt, using local prompt matching when the prompt is sufficiently close and an LLM retrieval agent otherwise. For the matched task, the training demonstrations provide a set of initial joint positions. Sci treats these initial states as samples from the feasible target-state distribution of the next atomic task, not as labels of the test composite trajectory. Therefore, retrieving $target_s$ does not use future states of the evaluated long-horizon sequence. It only selects one feasible starting state for the next atomic policy.
\begin{equation}
    \mathcal{P} = Extract(\mathcal{D})
\end{equation}
\begin{equation}
    target_s = Select(o^l_{i+1}, curr_s, \mathcal{P}, \mathcal{D})
\end{equation}

In principle, any initial state sampled from the next task's training demonstrations can be used as a target state, as long as the VLA can start the next atomic task from it. The retrieval objective is therefore not to reveal the test trajectory, but to choose a nearby and collision-free target state. In implementation, Sci ranks candidate target states by their joint-space distance to $curr_s$ and validates the top candidates by checking whether an interpolated path from $curr_s$ to the candidate introduces new collisions in simulation. This validation serves as a lightweight collision-prediction filter during retrieval: candidates whose interpolated restoration path intersects instruments, objects, or the table are rejected before the LLM generates transitional action code. The closest valid candidate is selected as $target_s$.

After obtaining the current observation, the next-task prompt, and $target_s$, a LLM coding agent generates transitional action code $\mathcal{C}_{i \rightarrow i+1}$.
\begin{equation}
    \mathcal{C}_{i \rightarrow i+1} = \texttt{LLM}(o^l_{i+1},o^v_t,target_s,curr_s)
\end{equation}

However, contents generated by LLMs may exhibit hallucination.
% {\color{blue}However, LLMs often suffer from semantic hallucinations.}
To reduce this risk and obtain executable target code, we use a two-stage prompt and a constrained code template. In the first stage, the LLM receives the front and side camera views, the coordinate-axis convention, and a safety-first objective, and returns a short JSON plan containing ordered steps such as releasing the gripper, translating the end effector along an axis, rotating the gripper, and moving away from visible obstacles. The prompt explicitly includes anti-collision heuristics: release the gripper before retreating when the end effector may still constrain an object, first move away from the nearest visible surface or instrument, prefer short axis-aligned translations through open space, avoid crossing visible objects or instrument lids, and rotate the gripper only when the rotation increases clearance. In the second stage, the LLM converts this plan into the body of the template's \texttt{execute} function using only allowed interfaces, including \texttt{get\_site\_pose}, \texttt{interpolate}, \texttt{path\_follow}, \texttt{move\_to}, \texttt{gripper\_control}, and \texttt{rotate\_gripper}. The template fixes the surrounding class structure, initializes IK outside the generated body, and inserts the final restoration to $target_s$ through \texttt{move\_to\_target\_qpos}. Thus, the LLM only specifies short obstacle-clearance motions, while target-state restoration is handled by the host code.
Once the restoration is successful, the VLA inference process can proceed. We first let the robotic arm avoid obstacles as much as possible, and then use the pose recovery function to reach the target state.
When the code is generated, control switches to the insertion module.

\textbf{Transitional Action Insertion Module.}
The insertion module is an action switch that toggles between the plugin and VLA.
We set the maximum runtime $T_i$ for each atomic task in advance
(in the experiment section, this is the average demonstration time per atomic task).
% {\color{blue}, which corresponds to the average demonstration duration.}
When the VLA reaches the final inference step, the module terminates the inference process regardless of the task completion status.
In practice, if we do not actively terminate the VLA inference, the robotic arm will continue to tremble.
% {\color{blue} this active termination prevents execution jittering in the robotic arm.}
Once the transitional action is generated by the generation module, VLA's inference is restored, allowing it to continue automatically executing the next atomic task.
% {\color{blue}Following the generation of transitional actions, SciVLA resumes VLA inference to automatically initiate the subsequent atomic task.}






\subsection{Digital Twin}

We utilize a simulated scientific laboratory system to evaluate our method.
This simulation facilitates efficient data collection and prevents potential damage to physical instruments.
In this paper, we employ the Autobio~\cite{lan2025autobio} laboratory environment, a biology lab built on the MuJoCo engine.
It incorporates essential instruments and equipment for biological research.
Experimental operations are performed using the UR5e robotic arm for long-horizon tasks.
Additionally, to provide a preliminary real-world demonstration, we
have constructed a
real-world scenario based on the simulation system to test the transferability of Sci from simulation to a real laboratory.



% \begin{figure}[!t]
%   \centering
%   \includegraphics[width=\linewidth]{Figures/new3dasset.pdf}
%   \caption{Some newly added instrument 3D assets examples. (a) Ozone cleaner. (b) Spin coater. (c) Hot plate. }
%   \label{figs:3d2}
% \end{figure}

\begin{figure*}[!t]
  \centering
  \includegraphics[width=\linewidth]{Figures/coherence.pdf}
  \caption{Action trajectory comparison for the cleaning-table sequence with the base $\pi_0$ model and Sci-$\pi_0$. The yellow lines represent the trajectory of the end effector. The dashed box indicates the trajectory of the transitional action.
  Simulated cleaning table task:
  pick 3 objects into the basket.
  }
  \label{figs:coherence}
\end{figure*}


\section{Experiments}
\label{sec:exp}
\subsection{Experiments Setup}

\textbf{Simulation Environment and Robot.}
All simulated experiments are conducted in the Autobio~\cite{lan2025autobio} laboratory environment, which is built on MuJoCo and contains common biological laboratory instruments. We use a UR5e robotic arm to execute all simulated manipulation sequences. The simulated environment is used for both demonstration collection and evaluation, allowing randomized object placements while avoiding damage to physical instruments during repeated long-horizon trials.

\textbf{VLA Baselines and Sci Variants.}
We evaluate Sci as an inference-time plugin for fine-tuned VLA models. For the ``cleaning table'' task, we select two diffusion-based VLAs, $\pi_0$~\cite{black2024pi_0} and $\pi_{0.5}$~\cite{black2025pi_}, and one autoregression-based VLA, $fast$~\cite{pertsch2025fast}, as base models. We denote the corresponding plugin-equipped models as Sci-$\pi_0$, Sci-$\pi_{0.5}$, and Sci-$fast$. For scientific operation tasks, we use $\pi_0$ as the base VLA and compare it with Sci-$\pi_0$. Sci inserts a transitional action between every two consecutive atomic tasks while keeping the VLA weights unchanged.

\textbf{Dataset Collection.}
All training demonstrations are collected at the atomic-task level rather than at the full long-task level. For the ``cleaning table'' task, we generate 100 demonstrations for each atomic task and fine-tune all three base models for 80,000 steps. During evaluation, the maximum inference time for each atomic task is set differently due to the complexity of the task. If the time limit is exceeded, inference for the current atomic task is terminated and the sequence immediately moves to the next atomic task. For scientific operations, we collect 100 demonstrations for each atomic task in simulation, covering 14 different atomic tasks in total. We use these collected atomic-task datasets to fine-tune $\pi_0$ for 50,000 steps on 2 NVIDIA H100 GPUs.

\textbf{Task Setup.}
First, we design a ``cleaning table'' task in pick-and-place scenarios. It is composed of 3 pick-and-place atomic tasks:
(1) \texttt{pick pipette tip box into the basket},
(2) \texttt{pick PCR plate into the basket}, and
(3) \texttt{pick centrifuge tube into the basket} (See in Fig.~\ref{figs:cleaning}).
The order of atomic tasks goes from simple to difficult.
For example, a PCR plate is harder to pick up than a pipette tip box because it is wider and thinner. Centrifuge tubes have a smaller diameter, requiring more precise handling.

Then, we design 6 long-horizon tasks specifically for scientific operations, with the visual layouts shown in Fig.~\ref{figs:taskABCDEF}. These long-horizon tasks are primarily designed around interactions between scientific instruments (e.g., centrifuges and thermal cyclers) and objects (e.g., tube racks, tubes, and PCR plates). The end state of the previous atomic task will affect the execution of the next atomic task. For example, if the centrifuge lid does not open, placing the centrifuge tubes will inevitably fail.
The atomic-task definitions used in the scientific operation experiments are summarized in Tab.~\ref{tab:long_task_definitions}.

\begin{table*}[!t]
\small
\begin{tabular}{cp{0.22\linewidth}p{0.65\linewidth}}
\toprule
Task & Long-horizon workflow & Atomic-task definitions \\
\midrule
A & Centrifuge loading and start &
A.1 opens the centrifuge 5910 lid $\rightarrow$ A.2 places the experimental centrifuge tube from the rack into the centrifuge $\rightarrow$ A.3 places the balance centrifuge tube from the rack into the centrifuge $\rightarrow$ A.4 closes the lid $\rightarrow$ A.5 presses the screen button to start the centrifuge. \\
\midrule
B & Centrifuge unloading-related operations &
B.1 opens the centrifuge lid $\rightarrow$ B.2 takes the experimental centrifuge tube from the centrifuge and places it on the rack $\rightarrow$ B.3 takes the balance centrifuge tube from the centrifuge and places it on the rack $\rightarrow$ B.4 closes the lid $\rightarrow$ B.5 presses the screen button. \\
\midrule
C & Thermal-cycler loading and start &
C.1 opens the thermal cycler lid $\rightarrow$ C.2 places the PCR plate into the thermal cycler $\rightarrow$ C.3 closes the lid $\rightarrow$ C.4 tightens the lid knob $\rightarrow$ C.5 presses the button to start the thermal cycler. \\
\midrule
D & Thermal-cycler unloading-related operations &
D.1 loosens the knob $\rightarrow$ D.2 opens the lid $\rightarrow$ D.3 takes the PCR plate from the thermal cycler $\rightarrow$ D.4 closes the lid $\rightarrow$ D.5 presses the button. \\
\midrule
E & Extended centrifuge workflow &
E.1 opens the centrifuge 5910 lid $\rightarrow$ E.2 places the experimental centrifuge tube from the rack into the centrifuge $\rightarrow$ E.3 places the balance centrifuge tube from the rack into the centrifuge $\rightarrow$ E.4 closes the lid $\rightarrow$ E.5 presses the screen button to start the centrifuge $\rightarrow$ E.6 opens the centrifuge lid again $\rightarrow$ E.7 takes the experimental centrifuge tube out $\rightarrow$ E.8 takes the balance centrifuge tube out. \\
\midrule
F & Extended thermal-cycler workflow &
F.1 opens the thermal cycler lid $\rightarrow$ F.2 places the PCR plate into the thermal cycler $\rightarrow$ F.3 closes the lid $\rightarrow$ F.4 tightens the lid knob $\rightarrow$ F.5 presses the button to start the thermal cycler $\rightarrow$ F.6 loosens the knob $\rightarrow$ F.7 opens the lid $\rightarrow$ F.8 takes the PCR plate from the thermal cycler. \\
\bottomrule
\end{tabular}
\caption{Definitions of the long-horizon scientific operation tasks. Each numbered entry (e.g., A.2 or F.7) corresponds to the same ``Atomic'' column in Tab.~\ref{tab:sci-ops}, so the table can be read without relying on the visual task diagram.}
\label{tab:long_task_definitions}
\end{table*}


We collect all these atomic tasks separately, and in each atomic task, we randomize the scenes. For example, test tubes are placed in a randomly selected slot on a rack, and PCR plates appear randomly within a designated area.
During inference, we fix the prompt for each long task and apply slight perturbations to the robotic arm's initial state.


\textbf{Evaluation Metrics.}
For each long task sequence, we let the base VLA model and the corresponding Sci-xxx model repeatedly test 20 times, where Sci-xxx denotes the same base model equipped with the Sci plugin. One inference refers to a long task from atomic task 1 to the last atomic task. We define from the end state of atomic task $i$ to the end state of atomic task $i+1$ as an atomic execution process. The success rate reported in the tables is the independent success rate of each atomic task: for each atomic-task column, we record whether that atomic execution process reaches its required final state over 20 complete sequence executions. Therefore, the table entries are per-atomic-task final-state success rates rather than full-sequence success rates.
Atomic 1 is evaluated from the long-task initial state, whereas Atomic $k$ ($k>1$) is evaluated after the preceding atomic task and, for Sci variants, after the inserted transitional action. A failure in a later atomic column can therefore arise either from an unsuitable transition state or from the intrinsic difficulty of the atomic manipulation policy. The results section distinguishes these cases when interpreting the table.
All the data in the tables is obtained from executing complete action sequences.


\begin{figure}[!t]
  \centering
  \includegraphics[width=\linewidth]{Figures/cleaning.pdf}
  \caption{The flow chart of ``cleaning table'' is composed of simple pick-and-place tasks. The order of execution goes from simple to difficult. Arrows represent the end effector's motion trajectory.}
  \label{figs:cleaning}
\end{figure}


\subsection{Cleaning Table Results}
\label{sec:basic}
The ``cleaning table'' task evaluates whether Sci can improve transitions in a simple pick-and-place sequence composed of three atomic tasks. This task is used as a basic validation because the three objects increase in manipulation difficulty while sharing the same basket-placement goal.

As shown in Tab.~\ref{tab:cleaning}, base VLA models only achieve good performance on the first atomic task.
Since our atomic task training data are collected separately, tasks can easily lack transitions.
The robotic arm enters an out-of-distribution state after finishing the first task.
This result indicates that the base VLA models can rely strongly on demonstration trajectories and can be sensitive to state shifts between atomic tasks.
In addition, the Sci-xxx models still have limited success rates on atomic task 2 and 3 when the base VLA model itself performs poorly. It can be seen that $\pi_0$ and $\pi_{0.5}$ generalize better than $fast$.
Compared to $\pi_0$, $\pi_{0.5}$ is better suited for more challenging atomic tasks.



\begin{table}[!t]
% \setlength{\tabcolsep}{1pt}
\begin{tabular}{cccc}
\toprule
\multirow{2}{*}{Method}&\multicolumn{3}{c}{Success Rate}\\
\cline{2-4}
&Atomic task 1&Atomic task 2&Atomic task 3\\
\midrule
$\pi_0$  & 14/20 (70\%) & 0/20 (0\%) & 0/20 (0\%)\\
\multirow{2}{*}{Sci-$\pi_0$} & \cellcolor{gray!25}14/20 (70\%) & \cellcolor{gray!25}\textbf{10/20 (50\%)} & \cellcolor{gray!25}4/20 (20\%)\\
&\cellcolor{gray!25}-&\cellcolor{gray!25}{\color{violet}$50\%\uparrow$}&\cellcolor{gray!25}{\color{violet}$20\%\uparrow$}\\ \hline

$\pi_{0.5}$  &13/20 (65\%) & 0/20 (0\%) & 0/20 (0\%) \\
\multirow{2}{*}{Sci-$\pi_{0.5}$} &\cellcolor{gray!25}12/20 (60\%)&\cellcolor{gray!25}3/20 (15\%)& \cellcolor{gray!25}\textbf{9/20 (45\%)}\\
&\cellcolor{gray!25}-&\cellcolor{gray!25}{\color{violet}$15\%\uparrow$}&\cellcolor{gray!25}{\color{violet}$45\%\uparrow$}\\ \hline

$fast$  &10/20 (50\%) & 0/20 (0\%) & 0/20 (0\%) \\
\multirow{2}{*}{Sci-$fast$} &\cellcolor{gray!25}10/20 (50\%) & \cellcolor{gray!25}5/20 (25\%) & \cellcolor{gray!25}5/20 (25\%) \\
&\cellcolor{gray!25}-&\cellcolor{gray!25}{\color{violet}$25\%\uparrow$}&\cellcolor{gray!25}{\color{violet}$25\%\uparrow$}\\
\bottomrule
\end{tabular}
\caption{Results of the baseline comparison for the ``cleaning table'' task. ``Sci-'' means the model uses the Sci inference plugin. Test a total of 20 times. Each test requires the model to complete an entire atomic task sequence. Data statistics report the independent final-state success rate of each atomic task.}
\label{tab:cleaning}
\end{table}


\subsubsection{Atomic Task Coherence}
In the ``cleaning table'' task, Sci-$\pi_0$ achieves higher task coherence than the base $\pi_0$ (shown in Fig.~\ref{figs:coherence}).
The base model $\pi_0$ is confused when starting the second atomic task, exhibiting continuous jitter. This jitter will not stop until the entire inference time is exhausted.
On the contrary, Sci-$\pi_0$ maintains a complete execution trajectory throughout all the atomic tasks.
Even if one atomic task fails, Sci-$\pi_0$'s transitional action can guide the robot toward a state from which subsequent tasks can start.



% \begin{figure*}[!ht]
%   \centering
%   \includegraphics[width=\linewidth]{Figures/taskABCD.pdf}
%   \caption{Results of four 5-element tasks in the centrifuge and thermal cycler operation scenario. Top: Schematic diagram of continuous atomic tasks for operating an instrument (task A, B, C and D in Tab.~\ref{tab:task}). The atomic tasks are related to each other and vary in difficulty. Bottom: Success rates of task inference repeated 20 times.}
%   \label{figs:taskABCD}
% \end{figure*}

% \begin{figure*}[!ht]
%   \centering
%   \includegraphics[width=\linewidth]{Figures/taskEF.pdf}
%   \caption{Results of two 8-element tasks in both two operation settings. Top: Schematic diagram of continuous atomic tasks for operation (task E and F in Tab.~\ref{tab:task}). The atomic tasks are related to each other and vary in difficulty. atomic tasks may appear repeatedly. Bottom: Success rates of task inference repeated 20 times.}
%   \label{figs:taskEF}
% \end{figure*}

\begin{figure*}[!t]
  \centering
  \includegraphics[width=\linewidth]{Figures/taskABCDEF.pdf}
  \caption{Results of scientific operations in the centrifuge and thermal cycler scenario. Top left: Flow charts of continuous atomic tasks for operating instruments. Top right: Images from the wrist camera of 3 atomic tasks with high precision requirements. Bottom: Success rates of each atomic task. Each sequence is tested 20 times.}
  \label{figs:taskABCDEF}
\end{figure*}


\subsection{Scientific Operation Results}
For scientific operations, the fine-tuned $\pi_0$ and Sci-$\pi_0$ are evaluated on six long-horizon task sequences under random initial disturbances. Tasks A--D contain five atomic tasks, while Tasks E--F extend the centrifuge and thermal-cycler workflows to eight atomic tasks. In Tab.~\ref{tab:sci-ops}, the columns ``Atomic 1''--``Atomic 8'' follow the numbered definitions in Tab.~\ref{tab:long_task_definitions}. For instance, Atomic 2 in Task C is C.2, placing the PCR plate into the thermal cycler.

% We still report the success rate for each atomic task, like in Sec.~\ref{sec:basic}.
As shown in Fig.~\ref{figs:taskABCDEF} and Tab.~\ref{tab:sci-ops}, Sci-$\pi_0$ improves the independent per-atomic-task success rates in scientific scenarios, including both 5-step and 8-step task sequences. Atomic 1 mainly evaluates whether the base VLA can start from the initial configuration, whereas each later atomic task requires both a suitable transition from the previous final state and successful execution of the next atomic policy. We therefore interpret gains on later non-placement operations, such as closing lids, tightening knobs, pressing buttons, or reopening instruments, as evidence that Sci reduces state-mismatch across atomic tasks. In contrast, low success on tube or PCR-plate placement and removal steps reflects a mixture of transition quality and the precision limits of the underlying atomic VLA policy.
{Notably}, we observe some counterintuitive phenomena in the baseline results: the high success rates for task A.4 (\texttt{close centrifuge lid}) and task C.3 (\texttt{close thermal cycler lid}).
The reason is that the initial states of task A.4 align with the end state of A.3. The gripper happened to be on the lid when performing task A.4. This indicates a strong reliance on the data among these base VLA models, as the limited training set does not include all joint positions.

\begin{figure}[!t]
  \centering
  \includegraphics[width=\linewidth]{Figures/pcr.pdf}
  \caption{Top: 3D assets of a thermal cycler slot (left) and a PCR plate (right).
  Bottom: corresponding real objects}
  \label{figs:pcr}
\end{figure}

These results also show a boundary of Sci: it focuses only on transitions between atomic tasks and does not address issues within those tasks. Therefore, for some tasks with high precision requirements, such as \texttt{placing centrifuge tube} or \texttt{placing a PCR plate into slots} (See top right in Fig.~\ref{figs:taskABCDEF}), the success rate depends largely on the VLA base model itself. Fig.~\ref{figs:pcr} is an example of the thermal cycler slot. A PCR plate usually has many wells, and placing it into a slot from above requires extremely high precision.
Additionally, PCR plates may be soft, and when the gripper holds the plate, deformation may occur. This situation significantly increases the difficulty of the atomic task.
Future research should prioritize improving operational precision and balancing the distribution of simple and complex demonstrations in datasets to mitigate training imbalances.


\begin{table*}[!ht]
    \setlength{\tabcolsep}{0.3pt}
  \begin{tabular}{cccccccccc}
    \toprule
    \multirow{2}{*}{Number}&\multirow{2}{*}{Method}&\multicolumn{8}{c}{Success Rate}\\
    \cline{3-10}
    &&Atomic 1&Atomic 2&Atomic 3&Atomic 4&Atomic 5&Atomic 6&Atomic 7&Atomic 8\\
    \midrule

    \multirow{3}{*}{A}& $\pi_0$ & 20/20 (100\%) & 0/20 (0\%) &0/20 (0\%)&12/20 (60\%)&0/20 (0\%)&-& -&-\\
    &\multirow{2}{*}{Sci-$\pi_0$} & \cellcolor{gray!25}20/20 (100\%) & \cellcolor{gray!25}\textbf{6/20 (30\%)} & \cellcolor{gray!25}\textbf{3/20 (15\%)} & \cellcolor{gray!25}\textbf{20/20 (100\%)} & \cellcolor{gray!25}\textbf{19/20 (95\%)} &\cellcolor{gray!25}-&\cellcolor{gray!25}-&\cellcolor{gray!25}-\\

    &&\cellcolor{gray!25}-&\cellcolor{gray!25}{\color{violet}$30\%\uparrow$}&\cellcolor{gray!25}{\color{violet}$15\%\uparrow$}&\cellcolor{gray!25}{\color{violet}$40\%\uparrow$}&\cellcolor{gray!25}{\color{violet}$95\%\uparrow$}&\cellcolor{gray!25}-&\cellcolor{gray!25}-&\cellcolor{gray!25}-\\

    \cline{1-10}
    \multirow{3}{*}{B}&$\pi_0$& 20/20 (100\%) & 0/20 (0\%) &0/20 (0\%)&0/20 (0\%)&0/20 (0\%)&-& -&-\\
    &\multirow{2}{*}{Sci-$\pi_0$} & \cellcolor{gray!25}20/20 (100\%) & \cellcolor{gray!25}\textbf{20/20 (100\%)} &\cellcolor{gray!25}\textbf{5/20 (25\%)}&\cellcolor{gray!25}0/20 (0\%)&\cellcolor{gray!25}\textbf{15/20 (75\%)}&\cellcolor{gray!25}-&\cellcolor{gray!25}-&\cellcolor{gray!25}-\\

    &&\cellcolor{gray!25}-&\cellcolor{gray!25}{\color{violet}$100\%\uparrow$}&\cellcolor{gray!25}{\color{violet}$25\%\uparrow$}&\cellcolor{gray!25}{-}&\cellcolor{gray!25}{\color{violet}$75\%\uparrow$}&\cellcolor{gray!25}-&\cellcolor{gray!25}-&\cellcolor{gray!25}-\\

    \cline{1-10}
    \multirow{3}{*}{C}&$\pi_0$& 18/20 (90\%) &0/20 (0\%) &12/20 (60\%)&0/20 (0\%)&0/20 (0\%)&-&-&-\\
    &\multirow{2}{*}{Sci-$\pi_0$} & \cellcolor{gray!25}20/20 (100\%) & \cellcolor{gray!25}\textbf{10/20 (50\%)} &\cellcolor{gray!25}\textbf{20/20 (100\%)}&\cellcolor{gray!25}\textbf{14/20 (70\%)}&\cellcolor{gray!25}\textbf{18/20 (90\%)}&\cellcolor{gray!25}-&\cellcolor{gray!25}-&\cellcolor{gray!25}-\\

    &&\cellcolor{gray!25}-&\cellcolor{gray!25}{\color{violet}$50\%\uparrow$}&\cellcolor{gray!25}{\color{violet}$40\%\uparrow$}&\cellcolor{gray!25}{\color{violet}$70\%\uparrow$}&\cellcolor{gray!25}{\color{violet}$90\%\uparrow$}&\cellcolor{gray!25}-&\cellcolor{gray!25}-&\cellcolor{gray!25}-\\

    \cline{1-10}
    \multirow{3}{*}{D}&$\pi_0$& 15/20 (75\%) & 0/20 (0\%) &0/20 (0\%)&0/20 (0\%)&0/20 (0\%)&-&-&-\\
    &\multirow{2}{*}{Sci-$\pi_0$} & \cellcolor{gray!25}16/20 (80\%) & \cellcolor{gray!25}\textbf{10/20 (50\%)} &\cellcolor{gray!25}\textbf{18/20 (90\%)}&\cellcolor{gray!25}\textbf{2/20 (10\%)}&\cellcolor{gray!25}\textbf{18/20 (90\%)}&\cellcolor{gray!25}-&\cellcolor{gray!25}-&\cellcolor{gray!25}-\\

    &&\cellcolor{gray!25}-&\cellcolor{gray!25}{\color{violet}$50\%\uparrow$}&\cellcolor{gray!25}{\color{violet}$90\%\uparrow$}&\cellcolor{gray!25}{\color{violet}$10\%\uparrow$}&\cellcolor{gray!25}{\color{violet}$90\%\uparrow$}&\cellcolor{gray!25}-&\cellcolor{gray!25}-&\cellcolor{gray!25}-\\

    \cline{1-10}
    \multirow{3}{*}{E}&$\pi_0$&  20/20 (100\%) & 0/20 (0\%)&0/20 (0\%)&10/20 (50\%)&0/20 (0\%)&0/20 (0\%)&0/20 (0\%)&0/20 (0\%)\\
    &\multirow{2}{*}{Sci-$\pi_0$} &  \cellcolor{gray!25}20/20 (100\%) & \cellcolor{gray!25}\textbf{3/20 (15\%)} & \cellcolor{gray!25}\textbf{2/20 (10\%)} & \cellcolor{gray!25}\textbf{20/20 (100\%)} & \cellcolor{gray!25}\textbf{17/20 (85\%)} &\cellcolor{gray!25}\textbf{15/20 (75\%)}&\cellcolor{gray!25}\textbf{2/20 (10\%)}&\cellcolor{gray!25}\textbf{13/20 (65\%)}\\

    &&\cellcolor{gray!25}-&\cellcolor{gray!25}{\color{violet}$15\%\uparrow$}&\cellcolor{gray!25}{\color{violet}$10\%\uparrow$}&\cellcolor{gray!25}{\color{violet}$50\%\uparrow$}&\cellcolor{gray!25}{\color{violet}$85\%\uparrow$}&\cellcolor{gray!25}{\color{violet}$75\%\uparrow$}&\cellcolor{gray!25}{\color{violet}$10\%\uparrow$}&\cellcolor{gray!25}{\color{violet}$65\%\uparrow$}\\

    \cline{1-10}
    \multirow{3}{*}{F}&$\pi_0$& 19/20 (95\%) & 0/20 (0\%)&15/20 (75\%)&0/20 (0\%)&0/20 (0\%)&0/20 (0\%)&0/20 (0\%)&0/20 (0\%)\\
    &\multirow{2}{*}{Sci-$\pi_0$} & \cellcolor{gray!25}19/20 (95\%) & \cellcolor{gray!25}\textbf{9/20 (45\%)} &\cellcolor{gray!25}\textbf{18/20 (90\%)}&\cellcolor{gray!25}\textbf{18/20 (90\%)}&\cellcolor{gray!25}\textbf{15/20 (75\%)}&\cellcolor{gray!25}\textbf{7/20 (35\%)}&\cellcolor{gray!25}\textbf{12/20 (60\%)}&\cellcolor{gray!25}0/20 (0\%)\\

    &&\cellcolor{gray!25}-&\cellcolor{gray!25}{\color{violet}$45\%\uparrow$}&\cellcolor{gray!25}{\color{violet}$15\%\uparrow$}&\cellcolor{gray!25}{\color{violet}$90\%\uparrow$}&\cellcolor{gray!25}{\color{violet}$75\%\uparrow$}&\cellcolor{gray!25}{\color{violet}$35\%\uparrow$}&\cellcolor{gray!25}{\color{violet}$60\%\uparrow$}&\cellcolor{gray!25}{-}\\
  \bottomrule
\end{tabular}
\caption{Independent per-atomic-task success rates when repeatedly performing multiple long-term science operation tasks 20 times. The meaning of each numbered atomic column is defined in Tab.~\ref{tab:long_task_definitions}. Entries therefore indicate which specific atomic step succeeded or failed rather than a full-sequence success rate.}
\label{tab:sci-ops}
\end{table*}

% \begin{table*}[!ht]
%     % \setlength{\tabcolsep}{2pt}
%   \caption{Results of repeatedly performing multiple long-term science operation tasks 20 times. }
%   \begin{tabular}{ccccccc}
%     \toprule
%     \multirow{2}{*}{Number}&\multirow{2}{*}{Method}&\multicolumn{5}{c}{Success Rate}\\
%     \cline{3-7}
%     &&atomic task 1&atomic task 2&atomic task 3&atomic task 4&atomic task 5\\
%     \midrule
%     \multirow{2}{*}{A}& $\pi_0$ & 20/20 (100\%) & 0/20 (0\%) &0/20 (0\%)&12/20 (60\%)&0/20 (0\%) \\
%     &$\pi_0$+AVIP & 20/20 (100\%) & \textbf{6/20 (30\%)} & \textbf{3/20 (15\%)} & \textbf{20/20 (100\%)} & \textbf{19/20 (95\%)} \\
%     \cline{1-7}
%     \multirow{2}{*}{B}&$\pi_0$ & 20/20 (100\%) & 0/20 (0\%) &0/20 (0\%)&0/20 (0\%)&0/20 (0\%) \\
%     &$\pi_0$+AVIP & 20/20 (100\%) & \textbf{20/20 (100\%)} &\textbf{5/20 (25\%)}&0/20 (0\%)&\textbf{15/20 (75\%)}\\
%     \cline{1-7}
%     \multirow{2}{*}{C}&$\pi_0$ & 18/20 (90\%) &0/20 (0\%) &12/20 (60\%)&0/20 (0\%)&0/20 (0\%) \\
%     &$\pi_0$+AVIP & 20/20 (100\%) & \textbf{10/20 (50\%)} &\textbf{20/20 (100\%)}&\textbf{14/20 (70\%)}&\textbf{18/20 (90\%)}\\
%     \cline{1-7}
%     \multirow{2}{*}{D}&$\pi_0$ & 15/20 (75\%) & 0/20 (0\%) &0/20 (0\%)&0/20 (0\%)&0/20 (0\%) \\
%     &$\pi_0$+AVIP & 16/20 (80\%) & \textbf{10/20 (50\%)} &\textbf{18/20 (90\%)}&\textbf{2/20 (10\%)}&\textbf{18/20 (90\%)}\\
%   \bottomrule
% \end{tabular}
% \end{table*}


% \begin{table*}[!ht]
% \setlength{\tabcolsep}{2pt}
%   \caption{Results of repeatedly performing multiple long-term science operation tasks 20 times. }
%   \begin{tabular}{cccccccccc}
%     \toprule
%     \multirow{2}{*}{Number}&\multirow{2}{*}{Method}&\multicolumn{8}{c}{Success Rate}\\
%     \cline{3-10}
%     &&atomic task 1&atomic task 2&atomic task 3&atomic task 4&atomic task 5&atomic task 6&atomic task 7&atomic task 8\\
%     \midrule
%     \multirow{2}{*}{E}&$\pi_0$ &  20/20 (100\%) & 0/20 (0\%)&0/20 (0\%)&10/20 (50\%)&0/20 (0\%)&0/20 (0\%)&0/20 (0\%)&0/20 (0\%) \\
%     &$\pi_0$+AVIP &  20/20 (100\%) & \textbf{3/20 (15\%)} & \textbf{2/20 (10\%)} & \textbf{20/20 (100\%)} & \textbf{17/20 (85\%)} &\textbf{15/20 (75\%)}&\textbf{2/20 (10\%)}&\textbf{13/20 (65\%)}\\
%     \cline{1-10}
%     \multirow{2}{*}{F}&$\pi_0$ & 19/20 (95\%) & 0/20 (0\%)&15/20 (75\%)&0/20 (0\%)&0/20 (0\%)&0/20 (0\%)&0/20 (0\%)&0/20 (0\%) \\
%     &$\pi_0$+AVIP & 19/20 (95\%) & \textbf{9/20 (45\%)} &\textbf{18/20 (90\%)}&\textbf{18/20 (90\%)}&\textbf{15/20 (75\%)}&\textbf{7/20 (35\%)}&\textbf{12/20 (60\%)}&0/20 (0\%)\\
%   \bottomrule
% \end{tabular}
% \end{table*}

\subsection{Ablation Studies}
We conduct ablation studies on the C.1$\rightarrow$C.2 transition because this sequence contains both an instrument-opening action and a subsequent high-precision placement action. Each setting is tested 10 times on the same long-task transition, and we report whether C.1 is completed, whether C.2 is completed after the transition, and whether the transition introduces a collision. The ablation evaluates the following scenarios:
\begin{enumerate}
    \item Direct interpolation utilizing current state and target state.
    \item Operation without the use of either current state or target state.
    \item Our proposed method: Sci.
\end{enumerate}

The results indicate that current qpos and target state are important inputs. The combination of target-state restoration and LLM-driven obstacle avoidance planning contributes to the success of the subsequent task.

\begin{table}[!t]
% \setlength{\tabcolsep}{1pt}
\begin{tabular}{lccc}
\toprule
Method & C.1 & C.2 & Collision \\
\midrule
LLM only & 9/10 & 1/10 & 0/10 \\
Direct interpolation & 9/10 & 0/10 & 10/10 \\
Ours & 9/10 & \textbf{7/10} & 0/10 \\
\bottomrule
\end{tabular}
\caption{Ablation results on the long task C.1$\rightarrow$C.2.}
\label{tab:ablation_long_task}
\end{table}

\subsection{Real-World Validation}
In this subsection, we provide a preliminary real-world validation of Sci in a laboratory setting. We construct a real-world task scenario based on Task C in Fig.~\ref{figs:taskABCDEF}: operating a thermal cycler. This composite task consists of two atomic tasks, C.1:\texttt{opening the thermal cycler lid} and C.2:\texttt{placing a PCR plate inside}. The validation uses the same transition setting as the simulated C.1$\rightarrow$C.2 task, where Sci generates a transitional action after C.1 and before C.2.

We fine-tune $\pi_0$ with these two atomic tasks. Fig.~\ref{figs:real} illustrates a live demonstration of this sequence, showing the decomposition of transitional action inferred by Sci. During this demonstration, the robotic arm first moves to a position away from the instrument cover. Then, the gripper releases, preparing the execution of the next atomic task. Finally, the robotic arm moves to the starting position of the next atomic task. This demonstration provides qualitative evidence that the simulated transition design can be executed in the physical laboratory setup.

\begin{figure*}[!t]
  \centering
  \includegraphics[width=\linewidth]{Figures/real_demo.pdf}
  \caption{Real scenario of operating the thermal cycler (Task $C1 \rightarrow C2$). The trajectories within the dashed box indicate transitional action.}
  \label{figs:real}
\end{figure*}

\subsection{Latency and Generation Quality}
The total execution time of the task includes the VLA inference time and the inference time of the Sci plugin. LLM inference time is dependent on model size and hardware specifications, as is VLA inference. Therefore, the proportion of the total task duration spent on LLM inference varies. To document this variability, we test different sizes of open-source models on the C.1$\rightarrow$C.2 task 10 times. We focus on open-source models in this comparison because access to proprietary models was not consistently available during evaluation.

In these tests, VLA inference is executed on a single NVIDIA H100 GPU, while local LLM inference uses a single NVIDIA A800 GPU. Larger models generate more reliable transitional action code in this setting, as reflected by the higher C.2 success count, but they also require longer inference time. Smaller local models reduce latency, but their generated transition code is less reliable for the high-precision C.2 placement task.

\begin{table}[t]
\begin{tabular}{lccc}
\toprule
Model & Average & Total & C.2 \\
\midrule
qwen3.5-397b-a17b & 185s (81\%) & 227s (100\%) & 8/10 \\
qwen3.5-35b-a3b (local) & 69s (60\%) & 115s (100\%) & 7/10 \\
qwen3.5-9b (local) & 55s (56\%) & 99s (100\%) & 5/10 \\
\bottomrule
\end{tabular}
\caption{Comparison of LLM inference time and task performance on the C.1$\rightarrow$C.2 task. ``Average'' reports the average Sci plugin inference time and its proportion of the total execution time.}
\label{tab:llm_inference_time}
\end{table}


\section{Conclusion}
This paper proposes Sci, an inference plugin for VLA models in scientific robot manipulation, requiring no VLA retraining. It aims to address the discontinuity in transitions between atomic tasks during high-intensity, continuous atomic-task operations at the VLA.
 Additionally, we augment the biological laboratory simulation environment provided by Autobio with objects and instruments commonly used in other laboratories (e.g., chemistry) and establish a test framework for continuous operational tasks in scientific laboratories.
 Experiments indicate that Sci can improve the independent success rate of subsequent atomic tasks and enhance their coherence.

Although our plugin includes some safety constraints in the code template,
it cannot completely guarantee the correctness of the action trajectories generated by the LLM.
In cases where collisions are likely, we will let LLM agents repeatedly generate new actions, which will slow down the inference of the entire task.
Future work could focus on aligning LLM-agent and VLA inference times to minimize latency. Or use a locally fine-tuned code agent for transitional action reasoning, which can eliminate the impact of network latency.

Meanwhile, our experiments indicate that VLA performance remains sensitive to the coverage and quality of atomic-task demonstrations. Future work should study how to collect more diverse transition-relevant demonstrations and how to improve VLA robustness when demonstration quality is uneven.

\bibliography{refs}

\input{ReproducibilityChecklist.tex}

\end{document}
